% Reconstructed from the compiled lecture-note PDF.
\chapter{The Reinhardt minimax problem}
\section{Convex disks and packings}

\begin{displaytext}
A convex disk is a compact convex subset K ⊂R2 with nonempty interior. The word “disk” is \
topological: K need not be round. A convex disk is centrally symmetric if there is a point c such \
that x ∈K implies 2c −x ∈K. Translating the set does not affect packing density, so we place the \
center at the origin and write simply \
K = −K.
\end{displaytext}

\begin{displaytext}
A family P of congruent copies of K is a packing if distinct copies have disjoint interiors. Let BR \
be the Euclidean ball of radius R centered at the origin. When the limit exists, the density of the \
packing is \
1 \
δ(K, P) = lim X area(K′ ∩BR). \
R→∞ area(BR) K′∈P
\end{displaytext}

The greatest packing density of K is


\begin{displaytext}
δ(K) = sup δ(K, P). \
P
\end{displaytext}

This is the best fraction of the plane that can be covered by pairwise nonoverlapping congruent copies of K.


\begin{insightbox}{Geometric intuition}
\end{insightbox}

There are two competing optimizations. For a fixed shape K, the packer tries to make δ(K) as large as possible. Reinhardt’s problem then asks the shape designer to make this best possible density as small as possible. Thus the objective is a minimax:


shape chooses K −→ packer chooses best arrangement .


A highly “unpackable” shape is one whose densest packing is still relatively sparse.


\section{Lattice packings}

A lattice in R2 is a discrete subgroup


L = Zv1 + Zv2


generated by linearly independent vectors v1, v2. Its determinant or covolume is


det(L) = |det(v1, v2)| .


\begin{displaytext}
The translates K + L = \{K + l: l∈L\} form a lattice packing when their interiors are disjoint. The \
density is then exactly \
area(K) \
δ(K, L) = . \
det(L)
\end{displaytext}

The reason is elementary: a fundamental parallelogram of L has area det(L) and contains exactly one translate of K modulo the lattice.


Define the greatest lattice-packing density by


\begin{displaytext}
δL(K) = sup \{δ(K, L) : K + L is a packing\} .
\end{displaytext}

A crucial two-dimensional theorem removes the apparent difference between arbitrary and lattice packings.


\begin{insightbox}{Standard theorem used as a black box}
\end{insightbox}

Fejes Tóth’s theorem. If K ⊂R2 is centrally symmetric and convex, then


\begin{displaytext}
δ(K) = δL(K).
\end{displaytext}

Thus the densest packing of such a disk can be realized by a lattice packing.


This theorem is special to the planar centrally symmetric setting. It is the first reason the Reinhardt problem is tractable.


\section{The minimax quantity}

Let Kccs denote all centrally symmetric convex disks centered at the origin. Set


\begin{displaytext}
δmin = inf δ(K). \
K∈Kccs
\end{displaytext}

The Reinhardt problem asks for both the numerical value of δmin and a classification of all minimizing shapes.


\begin{displaytext}
The disk has density \
π \
δ(B) = √ ≈0.906899. \
12 \
It is not the minimizer. Reinhardt found a slightly worse shape, the smoothed octagon, whose \
density is √ \
8 − 32 −log 2 \
δoct = √ ≈0.902414. \
8 −1
\end{displaytext}

\begin{figureplaceholder}
Figure 1.1: A schematic smoothed octagon. The dashed regular octagon is clipped at each vertex. The true rounded pieces are hyperbolic arcs; the drawing uses smooth surrogate arcs only for visualization.
\end{figureplaceholder}

\begin{statementbox}{Definition}
Definition 1.1 (Smoothed polygon). A smoothed polygon in the Reinhardt–Mahler sense is a centrally symmetric convex disk whose boundary is a finite cyclic concatenation of straight segments and arcs of hyperbolas. Each hyperbolic arc is tangent to the two adjacent straight segments, and its asymptotes are the lines extending the next pair of straight segments.
\end{statementbox}

The full conjecture is the following.


\begin{statementbox}{Theorem}
Theorem 1.2 (Reinhardt conjecture; open in the source paper). The smoothed octagon minimizes δ(K) over Kccs, and it is unique up to affine transformation.
\end{statementbox}

The proved theorem is weaker but structural.


\begin{statementbox}{Theorem}
Theorem 1.3 (Mahler’s First conjecture; proved in [HV]). Every global minimizer is a finite-sided smoothed polygon.
\end{statementbox}

\section{Affine invariance}

\begin{displaytext}
An invertible affine map has the form x 7→Ax + b with A ∈GL2(R). Translation is irrelevant, so \
consider A. Both the area of K and the covolume of a lattice are multiplied by |det A|. Hence
\end{displaytext}

\begin{displaytext}
area(AK) |det A| area(K) \
δ(AK, AL) = = = δ(K, L). \
det(AL) |det A| det(L)
\end{displaytext}

\begin{displaytext}
Taking suprema gives \
δ(AK) = δ(K).
\end{displaytext}

Thus the problem naturally lives modulo affine transformations. This freedom will later be used to normalize six special boundary points to the sixth roots of unity.


\section{Why a hexagon appears}

Suppose K + L is a lattice packing. A fundamental domain of the lattice can be chosen as a parallelogram, but one can often crop its corners and obtain a smaller centrally symmetric hexagon still containing K. Every centrally symmetric hexagon tiles the plane by translations. Consequently, the optimal lattice packing of a non-parallelogram K can be represented by a least-area centrally symmetric hexagon HK circumscribing K.


The exact theorem is proved in the next chapter. Its consequence is the formula


\begin{displaytext}
area(K) \
δ(K) = area(HK).
\end{displaytext}

\begin{displaytext}
Because the ratio is scale invariant, we can normalize \
√ \
area(HK) = 12.
\end{displaytext}

After this normalization, minimizing packing density is equivalent to minimizing area(K).


\begin{insightbox}{Proof roadmap}
\end{insightbox}

The original minimax problem now starts to collapse:


arbitrary packings of K


⇓Fejes Tóth


lattice packings


⇓critical hexagon


\begin{displaytext}
δ(K) = area(K)/ area(HK) \
√ \
⇓normalize area(HK) = 12
\end{displaytext}

minimize the area of K.


The remaining difficulty is that K still ranges over an infinite-dimensional space of convex curves.


\section{Historical orientation}

Blaschke recorded the problem in 1923 as a conjecture attributed to Courant, who expected the disk to be worst. Reinhardt disproved this in 1934 by constructing the smoothed octagon and developed the critical-hexagon geometry. Mahler independently arrived at the same candidate in the 1940s. In 1947 he wrote that the boundary of an extremal domain should consist of line segments and hyperbolic arcs, but he could not prove this. That structural prediction is what the source paper establishes.


\begin{insightbox}{Chapter summary}
\end{insightbox}

The Reinhardt problem is not to pack one fixed shape; it is to choose the centrally symmetric convex shape whose best packing is least dense. Fejes Tóth reduces arbitrary packings to lattice packings. A least-area centrally symmetric circumscribed hexagon then converts the density into an area ratio. After normalization, the objective is simply to minimize the area of a convex disk satisfying a rigid critical-hexagon condition.

